\documentclass[11pt]{article}
\usepackage{mathunal-base}
\mathunalfuente{serif}

\renewcommand{\examtitulo}{Cálculo en Varias Variables --- Tercer Parcial}
\renewcommand{\examfecha}{24 de agosto de 2019}

\begin{document}

\examheader

\vspace{4pt}
\noindent Nombre: \dotfill\ Cédula: \dotfill

\vspace{4pt}
\noindent Grupo: \dotfill\ Salón: \dotfill

\vspace{4pt}
\noindent Puntaje. Solo para uso oficial:\quad 1--4 \dotfill;\quad 5 \dotfill;\quad 6 \dotfill;\quad 7 \dotfill.\quad Total: \dotfill\quad Nota: \dotfill

\vspace{6pt}
\noindent\textbf{Instrucciones.} La duración del examen es de 1 hora y 45 minutos. El examen consta de siete preguntas en dos hojas impresas por ambos lados, en doble columna. Antes de empezar verifique que su examen esté completo y que su nombre y datos sean asignados a Usted. No desengrape su examen ni añada otras grapas. Utilice los espacios asignados para resolver cada pregunta. No está permitido: hacer preguntas, usar calculadoras, sacar hojas en blanco, ni tomar apuntes en hojas diferentes a las del examen. No se permite que ningún celular se esté a la vista, o sea usado.

\vspace{8pt}
\noindent\textit{En cada una de las preguntas 1 a 4 marque solamente todas las afirmaciones correctas cruzando con una ``x'' la caja correspondiente. Una marcación incorrecta anula una correcta. Se obtiene crédito de 6\,\% marcando todas las opciones correctas y no marcando ninguna de las incorrectas. Se obtiene crédito de 3\,\% con solo una opción mal escogida, y crédito de 0\,\% en todos los demás casos. En estas preguntas la calificación tendrá en cuenta sólo la respuesta.}

\begin{enumerate}
\item{} (6\,\%) Sea $C$ una curva paramétrica suave y orientada, y sea $C^-$ la misma curva pero orientada en sentido contrario. Sea $\vec F$ un campo vectorial, y $f$ una función escalar, ambos continuos en $\mathbb{R}^3$. Entonces:
\begin{itemize}
\item[] \fbox{\phantom{x}}\ (a) $\displaystyle \int_C f\,ds = -\int_{C^-} f\,ds$.
\item[] \fbox{\phantom{x}}\ (b) $\displaystyle \int_C \vec F \cdot d\vec r = \int_{C^-} \vec F \cdot d\vec r$.
\item[] \fbox{\phantom{x}}\ (c) $\displaystyle \int_C \vec F \cdot d\vec r = -\int_{C^-} \vec F \cdot d\vec r$.
\item[] \fbox{\phantom{x}}\ (d) $\displaystyle \int_C f\,ds = \int_{C^-} f\,ds$.
\end{itemize}

\item{} (6\,\%) Marque las afirmaciones correctas:
\begin{itemize}
\item[] \fbox{\phantom{x}}\ (a) Si $\vec F(x,y,z) = (x^3+e^x,\ -z^2+y^3,\ z^3)$, y $H = \{(x,y,z)\in\mathbb{R}^3 \mid x^2+y^2+z^2 = 4,\ z \geq 0\}$ y $P = \{(x,y,z)\in\mathbb{R}^3 \mid z = 4-x^2-y^2,\ z \geq 0\}$, entonces $\displaystyle \iint_H (\vec\nabla \times \vec F)\cdot d\vec s = \iint_P (\vec\nabla \times \vec F)\cdot d\vec s$.
\item[] \fbox{\phantom{x}}\ (b) Si $f$ tiene derivadas parciales continuas en $\mathbb{R}^3$, y $C$ es cualquier circunferencia, entonces $\displaystyle \int_C (\vec\nabla f)\cdot d\vec r = 0$.
\item[] \fbox{\phantom{x}}\ (c) $\vec F(x,y,z) = (xy,\ yz,\ zx)$ es un campo irrotacional, es decir $\operatorname{rot}\vec F = \vec 0$.
\item[] \fbox{\phantom{x}}\ (d) Existe un campo vectorial $\vec G$ en $\mathbb{R}^3$ tal que $\operatorname{rot}\vec G = (x\sin y,\ \cos y,\ z - xy)$.
\end{itemize}

\item{} (6\,\%) Sea $E$ la esfera sólida $x^2+y^2+z^2 \leq 3$. Si $S = \partial E$ y $\vec F(x,y,z) = 3x\,\vec i + z^3\,\vec j + 3y^3\,\vec k$, entonces:
\begin{itemize}
\item[] \fbox{\phantom{x}}\ (a) $\displaystyle \iint_S \vec F \cdot d\vec s = 108\pi$.
\item[] \fbox{\phantom{x}}\ (b) $\displaystyle \iint_S \tfrac13\,dS = 12\pi$.
\item[] \fbox{\phantom{x}}\ (c) $\displaystyle \iint_S 3\,dS = 36\pi$.
\item[] \fbox{\phantom{x}}\ (d) $\displaystyle \iint_S \vec F \cdot d\vec s = 12\sqrt3\,\pi$.
\end{itemize}

\item{} (6\,\%) Sea $\vec F(x,y) = (P(x,y),\ Q(x,y))$ campo conservativo con $\vec F = \vec\nabla f$, con $f$ de clase $C^2$ y $C$ es una parábola dada por $\vec r(x) = (x, x^2)$ con $-1 \leq x \leq 1$, entonces:
\begin{itemize}
\item[] \fbox{\phantom{x}}\ (a) $\dfrac{\partial P}{\partial y} - \dfrac{\partial Q}{\partial x} = 0$.
\item[] \fbox{\phantom{x}}\ (b) $\displaystyle \int_C \vec F \cdot d\vec r = f(-1,1) - f(1,1)$.
\item[] \fbox{\phantom{x}}\ (c) $\displaystyle \int_C \vec F \cdot d\vec r = f(1,1) - f(-1,1)$.
\item[] \fbox{\phantom{x}}\ (d) $\dfrac{\partial P}{\partial x} - \dfrac{\partial Q}{\partial y} = 0$.
\end{itemize}
\end{enumerate}

\vspace{4pt}
\noindent\textsc{Preguntas con procedimiento.} En las preguntas 5 a 7 responda mostrando \textit{detalladamente} el procedimiento de cálculo empleado. Indique los teoremas o procesos que sustentan sus cálculos.

\begin{enumerate}[start=5]
\item{} (25\,\%) Suponga que un campo vectorial está dado por $\vec F(x,y,z) = (-x,\ -y,\ 1)$. Calcule el flujo de $\vec F$ a través de la superficie $S$ con
\[
S = \left\{(x,y,z)\ \middle|\ z = \sqrt{x^2+y^2},\ 1 \leq x^2+y^2 \leq 4\right\},
\]
la cual está orientada con la normal apuntando hacia arriba. Incluya bosquejo de las curvas, áreas o volúmenes de integración involucrados. Indique, si lo hay, la parametrización y proyección $D$ usada en cada integral.

\item{} (25\,\%) Evalúe las integrales de línea:
\begin{itemize}
\item[(a)] (10\,\%) $\displaystyle \int_C y\cos(z)\,ds$, donde $C$ es la porción de hélice circular dada por $x = \cos(t)$, $y = \operatorname{sen}(t)$, $z = t$, con $0 \leq t \leq \pi/2$.
\item[(b)] (15\,\%) $\displaystyle \int_{C'} \vec F \cdot d\vec r$, con $\vec F = -x^4y\,\vec i + \left(x^2 + e^{\operatorname{sen}(y)}\right)\vec j$ y $C'$ es el cuadrado orientado positivamente (contrahorario) con vértices $(0,0)$, $(1,0)$, $(1,1)$ y $(0,1)$.
\end{itemize}

\item{} (26\,\%) Encuentre el trabajo realizado por el campo de fuerzas $\vec F = xy\,\vec i + xy\,\vec j + xz\,\vec k$ al mover, en sentido antihorario visto desde $(0, 10, 10)$, una partícula a lo largo de la curva de intersección del plano $y + z = 2$ y el cilindro $x^2 + y^2 = 9$. Incluya bosquejo de las curvas, áreas o volúmenes de integración involucrados. Indique, si lo hay, la parametrización y proyección $D$ usada en cada integral.
\end{enumerate}

\examfooter

\end{document}
